%!TEX root = ./../../springer-mono-format/book.tex
%!TEX root = ./../../thesis-trier.tex

Here, in \cref{syncmax:thm:Aprim-idempt} we characterize $A$-primitivity in transitive groups 
by synchronization of all one-point singulars.
We then note a few consequences. At the end, we also introduce the more general class
of unimodular transformations and characterize strongly $\{a,b\}$-primitiveness
by synchronization of transformations from this class. This is in fact the most general class
of transformations synchronized by strongly $\{a,b\}$-primitive groups, but 
we cannot generalize this to transitive $A$-primitive group for sets $A$ with more than three points.

\label{def:quasiorderidemp}
We use the following quasiorder in the present section:
for a transformation monoid  $T \le_M \mathcal T_n$,
write $a \rightsquigarrow_T b$ 
if $a = b$ or there exists an idempotent one-point singular $t \in T$
such that $t(a) = b$ (note that for $a\ne b$ this implies $\dupl{t} = \{b\}$).
\footnote{Two distinct points $a,b \in [n]$ are also in relation for the following types of 
	transformations $t \colon [n] \to [n]$:
	\begin{enumerate}
		\item if $t$ is a one-point singular that permutes its image with $t(a) = t(b)$ and $t^p(b) = b$ for some $p > 0$, or
		
		\item $t$ is a one-point singular with $t^j(a) = b$ for some $j \ge 1$ and $t(b) = b$.
	\end{enumerate}	
	This follows since then we can find  $i > 0$ such that $t^i$ is an idempotent one-point singular
	with $t^i(a) = b$.
	By \cref{prelim:lem:finite-semigrp-idempotent-power}
	we find $i > 0$ such that $t^i$ is idempotent.
	Then, for (1), we have $\dupl{t_{|t([n])}} = \emptyset$
	and, with \cref{prelim:eqn:dupl0}, this yields that $t^i$ is a one-point singular
	with image $t([n])$. Then $t^i(a) = t^i(b)$
	and since $b \in t([n])$,
	we have $t^i(b) = b$.
	For (2), we have $\dupl{t} = \{b\}$ and as in the proof of \cref{prelim:lem:transitive-grp-one-point-singular} 
	we can conclude that $t^i$ is an idempotent one-point singular.
	We then must have $t^i(a) = b$
	since otherwise $i < j$ and $t^i(a)$ is not contained in a cycle, contradicting
	$t^i(t^i(a)) = t^i(a)$.
	%	
	% Umkehrung: 
	%	Next, suppose $t^i$ is an idempotent one-point singular with $t(a) = b$.
	%	Then~$t$ must be a one-point singular.
	%	%Assume $t(b) \notin \{b, t(a)\}$.
	%	%Assume $t(b) \ne b$.
	%	
	%	Suppose $t$ does not permute its image.
	%	Then $t(t(c))$ for some $c \in [n]$
	%	ends at the duplicate state and $t(c)$ and~$c$ are not contained in any cycle.
	%	If the cycle is not a self-loop, then $t(t(c))$ and $t(c)$ are mapped to distinct states on the cycle
	%	and $t^2$ is not a one-point singular. So, this is not possible.
	%	Thus, the cycle that contains the duplicate state is a self-loop	
}
If the transformation
monoid $T$ is clear from context, we shall simply write~$\rightsquigarrow$ instead.

\begin{lemma}
	\label{sync_max:lem:idempotent_quasiorder}
	Let $T \le_M \mathcal T_n$ be a transformation monoid with group of units $G$.
	Then~$\rightsquigarrow$ is a quasiorder
	such that, for $a,b \in [n]$ and $g \in G$,
	$a \rightsquigarrow b$
	implies $\applyfunc{a}{g} \rightsquigarrow \applyfunc{b}{g}$.
\end{lemma}
\begin{proof}
	By definition, $\rightsquigarrow$ is reflexive.
	Next, let $a,b,c \in [n]$
	be distinct such that $a \rightsquigarrow b$
	and $b \rightsquigarrow c$,	witnessed by idempotent one-point singulars
	$h_1, h_2 \in T$ 
	such that $h_1(a) = b$
	and $h_2(b) = c$.
	The points $b$ and $c$ are the unique duplicate points of $h_1$ and~$h_2$, respectively.
	Thus $\dupl{h_1} \subseteq h_1([n]) \cap h_2^{-1}(\dupl{h_2})$
	and so, by \cref{prelim:trans:idempt-one-point:prop:comp}
	the composition $\concat{h_1}{h_2}$
	is an idempotent one-point singular.
	We have $h_2(h_1(a)) = h_2(b) = c$.
	We also have $h_1(c) = c$
	or $h_1(c) = b$ and both cases give $h_2(h_1(c)) = c$.
	Thus we have $(\concat{h_1}{h_2})(a) = c$
	and so $a \rightsquigarrow c$. 
	Next, let $g \in G$.
	By \cref{prelim:conj}
	we have $\applyfunc{\applyfunc{a}{g}}{h_1^g} = \applyfunc{b}{g}$
	and $h_1^g$ is an idempotent
	one-point singular.
	So we find $\applyfunc{a}{g} \rightsquigarrow  \applyfunc{b}{g}$.
\end{proof}

\begin{theorem} \label{syncmax:thm:Aprim-idempt}
Suppose $G \le_G \mathcal S_n$ is transitive and $A \subseteq [n]$ is a subset with $|A| \ge 2$. Then the following are equivalent. 
\begin{enumerate}
\item \label{syncmax:thm:Aprim-idempt:1} $G$ is $A$-primitive.

\item \label{syncmax:thm:Aprim-idempt:2} $G$ synchronizes every one-point singular $t$ on~$[n]$
with $A \subseteq t^{-1}(\dupl{t})$. 

\item \label{syncmax:thm:Aprim-idempt:3} $G$ synchronizes some idempotent one-point singular $t$ on~$[n]$
with $t^{-1}(\dupl{t}) = A$.
\end{enumerate}
\end{theorem}
\begin{proof}
\ref{syncmax:thm:Aprim-idempt:1} $\Rightarrow$ \ref{syncmax:thm:Aprim-idempt:2}:
Since~$G$ is transitive, by \cref{prelim:lem:transitive-grp-one-point-singular}, there
exists an idempotent one-point singular $h \in \langle G \cup \{t\} \rangle_S$
with $A \subseteq h^{-1}(\dupl{h})$.
Let $A = \{a_1, \ldots, a_m\}$ for $m \ge 2$
and $h(a_i) = a_m$ for $i \in \{1,\ldots,m\}$.
Pick $i \in \{1,\ldots,m\}$.  We have $a_i \rightsquigarrow a_m$ since $h(a_i) = a_m$.
Denote by~$\leftrightsquigarrow$
the equivalence relation
induced by the quasiorder~$\rightsquigarrow$,
which is a $G$-congruence
by \cref{sync_max:lem:idempotent_quasiorder}.
Since~$G$ is transitive, there exists $g \in G$
with $\applyfunc{a_1}{g} = a_m$.
Pick some $k > 0$ 
such that $\applyfunc{a_i}{g^k} = a_i$.
By successively applying~$g$
to $a_i \rightsquigarrow a_m$
we find $a_i \rightsquigarrow \applyfunc{a_i}{g}, 
\applyfunc{a_i}{g} \rightsquigarrow \applyfunc{a_i}{g^2},
\ldots, 
\applyfunc{a_i}{g^{k-1}} \rightsquigarrow a_i$.
So $a_i \leftrightsquigarrow \applyfunc{a_i}{g}$ where $\applyfunc{a_i}{g} = a_m$.
Thus, we have $a_i \leftrightsquigarrow a_m$ for all $i \in \{1,\ldots,m\}$.
Since~$G$ is $A$-primitive,
this implies that all
pairs of states are equivalent.
Hence, for all pairs $c,d \in [n]$
there exists an idempotent $f \in \langle G \cup \{t\} \rangle_S$
such that $\applyfunc{c}{f} = \applyfunc{d}{f} = d$.
So, with the transformation semigroup version 
of the pair criterion (\cref{prelim:remarks_trans_semigrp_semiautomata} and \cref{prelim:thm:pair_sync_condition}) it follows
that $\langle G \cup \{t\} \rangle_S$
contains a constant map.	

\ref{syncmax:thm:Aprim-idempt:2} $\Rightarrow$ \ref{syncmax:thm:Aprim-idempt:3}: This is clear.

\ref{syncmax:thm:Aprim-idempt:3} $\Rightarrow$ \ref{syncmax:thm:Aprim-idempt:1}: 
 If~$G$ is not $A$-primitive, then, by \cref{sync_max:cor:syncSI-abprim1}, it synchronizes no
 idempotent $t \colon [n] \to [n]$ with $t^{-1}(\dupl{t}) \subseteq A$. This gives the claim.
\end{proof}

\begin{remark}
	We note that instead of quasiorders in the proof of 
	the implication ``$\ref{syncmax:thm:Aprim-idempt:1} \Rightarrow \ref{syncmax:thm:Aprim-idempt:2}$''
	in \cref{syncmax:thm:Aprim-idempt},
	we can also construct
	a digraph on $[n]$ 
	that has an edge $(c,d)$ for each idempotent one-point singular
	mapping $c$ to $d$.
	By \cref{prelim:conj}, $G$ is a subgroup of the automorphism group
	of this graph. Furthermore, as~$G$ is transitive, this graph is strongly connected by \cref{prelim:thm:vertex_transitive_finite_graph}, 
	and paths correspond
	to idempotent one-point singulars mapping the starting vertex
	to the end vertex. Then, all pairs of points can be mapped to a common point
	by following a path, 
	which yields claim.
\end{remark}

\begin{corollary} \label{syncmax:cor:Aprim-idempt} \label{sync_max:thm:idemp}
Let $G \le_G \mathcal S_n$, $n \ge 3$ and $a,b \in [n]$ be distinct. Then the following are equivalent.
\begin{enumerate}
\item \label{syncmax:cor:Aprim-idempt:1} $G$ is strongly $\{a,b\}$-primitive.

\item \label{syncmax:cor:Aprim-idempt:2} $G$ synchronizes every one-point singular $t \colon [n] \to [n]$ with
$t(a) = t(b)$.

\item \label{syncmax:cor:Aprim-idempt:4} $G$ synchronizes every simple idempotent $t \colon [n] \to [n]$ with $t(a) = t(b)$. 
\end{enumerate}
\end{corollary}
\begin{proof}
\ref{syncmax:cor:Aprim-idempt:1} $\Rightarrow$ \ref{syncmax:cor:Aprim-idempt:2}:
 This follows from \cref{syncmax:thm:Aprim-idempt}.

\ref{syncmax:cor:Aprim-idempt:2} $\Rightarrow$ \ref{syncmax:cor:Aprim-idempt:4}: This is clear.

\ref{syncmax:cor:Aprim-idempt:4} $\Rightarrow$ \ref{syncmax:cor:Aprim-idempt:1}: This follows from \cref{sync_max:cor:2-sets_transitive}.
\end{proof}

\begin{remark}
 \Cref{syncaut:permgrp:ex:threesets} gives a transitive permutation group that
 is not $\{a,b\}$-primitive for all pairs of points $a,b \in [n]$,
 and so does not synchronize some simple idempotent,
 but which is $A$-primitive for all $A \subseteq [n]$
 with $|A| \ge 3$.
 Thus, synchronization of all idempotent one-point singulars
 with $|t^{-1}(\dupl{t})| = 3$
 is not sufficient for $\{a,b\}$-primitivity for some $\{a,b\}$.	
\end{remark}

As the trivial subgroup on $[2]$, which is not strongly $\{0,1\}$-primitive,
synchronizes every non-permutation,
\cref{syncmax:cor:Aprim-idempt} does not hold for $n = 2$.

By~\cref{sync_max:cor:syncSI-abprim2}, 
if $G$ synchronizes one simple idempotent $t$
with $t(a) = t(b)$, 
then $G$ is $\{a,b\}$-primitive.
However, as \cref{sync_max:ex:abprim-no-constant}
shows, $\{a,b\}$-primitivity
does not imply that $G$
synchronizes some simple idempotent $t$
with $t(a) = t(b)$.
\Cref{sync_max:thm:idemp} 
states that synchronization of both simple idempotents
mapping $a$ and $b$ to the same
point is equivalent to strong $\{a,b\}$-primitivity.

With our results we obtain a characterization of primitivity first given by 
Ara\'ujo and Cameron~\cite[Thm.~2]{AraujoCameron2014}.

\begin{corollary}[{Ara\'ujo and Cameron~\cite[Thm.~2]{AraujoCameron2014}}] \label{sync_max:thm:primtiive_kerneltype_k1111}
 Let $G \le_G \mathcal S_n$ be transitive and $1 < k < n$. Then~$G$ admits a non-trivial block of size at least~$k$
 if and only if there exists a one-point singular of defect $k - 1$ not synchronized by~$G$.
\end{corollary}
\begin{proof}
 First, suppose~$G$ has a non-trivial block $\Delta$ of size at least~$k$.
 Thus, $G$ is not $\Delta$-primitive.
 By \crefitem{syncmax:thm:Aprim-idempt}{syncmax:thm:Aprim-idempt:3}, 
 no idempotent one-point singular $t \colon [n] \to [n]$
 with $t^{-1}(\dupl{t}) = \Delta$
 is synchronized by~$G$
 and such a $t$ has defect $k - 1$.
 Conversely, suppose~$G$ does not synchronize some one-point singular $t \colon [n] \to [n]$ 
 of defect $k - 1$. By \cref{syncmax:thm:Aprim-idempt},
 $G$ is not $A$-primitive for $A = t^{-1}(\dupl{t})$
 and so there exists a non-trivial block
 of size at least $|A|$. By \cref{prop:char-one-attractor}, we have $|A| = k$.
\end{proof}

With \cref{sync_max:prop:prim_stronglyabprim,sync_max:thm:idemp} we also easily obtain
that a permutation group on $[n]$ where $n \ge 3$ is primitive
if and only if it synchronizes every idempotent one-point singular.

Since in a transitive group $G \le_G \mathcal S_n$
a block $\Delta$ and its equal-sized translates $g(\Delta)$
partition the domain, $n$ is a multiple of the size of every block
in a transitive permutation group. Thus \cref{sync_max:thm:primtiive_kerneltype_k1111} 
implies that a transitive
permutation group $G \le_G \mathcal S_n$ 
synchronizes every transformation
of kernel type $(k,1,\ldots,1)$
with $\frac{n}{2} < k < n$,
or that an imprimitive permutation group
with non-trivial blocks of size at most~$l$
synchronizes every transformation of kernel type $(k,1,\ldots,1)$
for $l < k < n$.
Example~\ref{syncaut:permgrp:ex:threesets}
gave imprimitive groups of arbitrarily large size that synchronize
every transformation of kernel type $(k,1,\ldots,1)$ with $k \ge 3$
as these groups only have non-trivial blocks of size two.

Lastly, we give a more general synchronization property for strongly $\{a,b\}$-primitive groups
and so-called unimodular transformations introduced in~\cite{Volkov24}.
This is based on two lemmas from~\cite[Lem.~4.3 \& 4.4]{KariVolkov2021} 
originally used for the solution of the road coloring problem in~\cite{KariVolkov2021}.
For the convenience of the reader we give full proofs (our lemmas are actually a little less general
than the original ones from~\cite{KariVolkov2021} and adapted to our setting of groups and non-permutations).
We will use the stability relation $\sim_s$ introduced on page~\pageref{def:stability}.

Let $t \colon [n] \to [n]$. For $a \in [n]$, set $l(a) = \min\{ i \mid \exists p > 0 \colon t^i(a) = t^{i + p}(a) \}$, i.e., the
shortest distance to a cycle state for the transformation induced by~$a$
(which is the shortest distance to a root state for a tail state).
Let $L = \max\{ l(a) \mid a \in [n] \}$.
We call $t$ \emph{unimodal}~\cite{Volkov24} if $L \ge 1$ and for $a,b \in [n]$ with $l(a) = l(b) = L$
we have $t^L(a) = t^L(b)$. In this case, two points $a,b \in [n]$
such that
\begin{enumerate}
	\item $t(b) = t^L(c)$ and $a = t^{L -1}(c)$ for some $c \in [n]$ with $l(c) = L$, and
	
	\item $t^p(b) = b$ for some $p > 0$, i.e., $b$ lies in a cycle of the functional graph of~$t$,
\end{enumerate} 
are called a \emph{unimodal pair} for~$t$.

Intuitively, a unimodular transformation is a non-permutation
with the property that all paths of maximal length that traverse a single cycle state (which must be an end state)
share the same cycle state as their end state. See \cref{fig:unimodal} for an example (taken from~\cite{Volkov24}).
One-point singulars are unimodal.

	\begin{figure}[htb]
	\centering
	
	\begin{tikzpicture}[
		scale=1.0,
		>=Stealth,
		every node/.style={circle, draw, inner sep=2pt},
		]			
		
		\node (A)  at (-2,2) {2};
		\node (B)  at (0,2)   {1};
		\node (C)  at (0,0.1)    {2};
		\node (D)  at (0,-1.8)  {1};
		\node (E)  at (4,-1.8) {1};
		
		\node (Y)  at (10,2) {3};
		
		\node (AD) at (2,0.1)   {0};
		\node (AB) at (2,2)  {0};
		\node (BC) at (4,2)  {0};
		
		\node (CD) at (4,0.1)   {0};
		\node (AC) at (6,2)  {1};
		\node (BD) at (8,2)  {2};
		
		\node (AE) at (6,0.1)   {2};
		\node (BE) at (9,0.1)   {3};
		
		\node (X)  at (8,-1.8) {3};
		
		
		\draw[->] (CD) -- (AD);	
		
		\draw[->] (AD) -- (AB);
		\draw[->] (AB) -- (BC);
		\draw[->] (BC) -- (CD);
		
		\draw[->] (A) -- (B);
		\draw[->] (C) -- (B);
		
		\draw[->] (D) --  (AD);
		\draw[->] (E) --  (AD);
		
		\draw[->] (B) -- (AB);
		
		\draw[->] (AC) -- (BC);
		
		\draw[->] (BD) -- (AC);
		
		\draw[->] (AE) -- (AC);
		
		\draw[->] (BE) --  (AE);
		
		\draw[->] (Y) -- (BD);
		
		\draw[->] (X) -- (AE);
		
		
		\draw[dashed,-] (A) -- (C);			
		\draw[dashed,-] (D) -- (E);
		\draw[dashed,-] (E) -- (CD);			
		\draw[dashed,-] (X) -- (BE);			
		\draw[dashed,-] (D) -- (CD);			
		\draw[dashed,-] (B) -- (AD);			
		\draw[dashed,-] (BD) -- (AE);
		
		\draw[dashed,-] (AB) to[out=45,in=135] (AC);
		
	\end{tikzpicture}
	\caption[Unimodal transformation]{Functional graph of a unimodal transformation.
		Marks indicate $l(a)$; we have $L = 3$. Dashed lines show the kernel relation.}
	\label{fig:unimodal}
\end{figure}

\IN{incompressible set}
For $S \le_S \mathcal T_X$ we call $A \subseteq X$ \emph{incompressible}
if there exists no $t \in S$ such that $|A| < |t(A)|$. Clearly, if~$A$
is incompressible, then $t(A)$ is incompressible for each $t \in S$
and $|A| = |t(A)|$.

\begin{lemma} \label{syncmax:lem:incompressible-stable-pair}
	Let $S \le_S \mathcal T_n$ be a semigroup and $A \subseteq [n]$
	be a maximal incompressible set.
	Suppose that there exists $t \in S$ that fixes every point in~$A$
	except one point $a \in A$. Then $a \sim_s t(a)$.
\end{lemma}
\begin{proof}
	Let $h \in S$. To show $a \sim_s t(a)$, we have to show that $f(h(a)) = f(h(t(a)))$
	for some $f \in S$.
	We may assume $h(a) \ne h(t(a))$.
	Set $B = A \setminus \{a\}$. Then we have the following:	
	\begin{enumerate}
		\item The set $h(t(A))$ is incompressible and has the same size as~$A$
		
		\item The set $h(A)$ is incompressible and has the same size as~$A$.
		
		\item $h(t(B)) = h(B)$.
	\end{enumerate}
	This yields that the set $C = h(B) \cup \{ h(a), h(t(a)) \}$
	has size $|A| + 1$.	
	Then, by the maximality of~$A$, there exists a pair of points in $C$ that can be mapped to the same point. Since $h(t(A))$ and $h(A)$ are incompressible, this pair must be $h(a)$ and $h(t(a))$.	
	Thus, we have $f(h(a)) = f(h(t(a)))$ for some $f \in S$.
	This yields $a \sim_s t(a)$.
\end{proof}

\begin{lemma} \label{syncmax:lem:unimodal-stable}
	Let $G \le_G \mathcal S_n$ be transitive
	and $t \colon [n] \to [n]$ be unimodal 
	with unimodal pair $a,b \in [n]$.
	Then $a \sim_s b$ w.r.t.~the transformation semigroup $\langle G \cup \{t\} \rangle_S$.
\end{lemma}
\begin{proof}
	Let $A = \{ c \in [n] \mid l(c) = L \}$.
	Since~$a, b \in [n]$ is a unimodal pair, without loss of generality,
	we have $t(a) = t(b) = t^L(c)$ and $a = t^{L-1}(c)$
	for some $c \in A$ and $b$ is contained in a cycle
	of the functional graph of~$t$.
	
	Next, consider a maximal incompressible $B \subseteq [n]$
	w.r.t.~$\langle G \cup \{t\} \rangle$.
	Let $d \in B$.
	By transitivity, we have $g(d) = c$ for some $g \in G$.
	The set $g(B)$ is also incompressbile.
	Thus, since $t^L(d) = t^L(d')$ for all $d,d' \in A$,
	there can only be one state from~$A$ in $g(B)$
	and so we have $g(B) \cap A = \{ c \}$.
	Let $C = t^{L - 1}(g(B))$.
	The set $C$ is incompressible and, by choice of~$L$,
	all of its points except $a = t^{L - 1}(c)$
	are contained in a cycle of the functional graph of~$t$.
	Let $m$ be the least common multiple of the lenghts
	of the cycles in the functional graph of~$t$.
	Then $t^m$ fixes all states in the cycles
	but, since $t^{m + 1} = t^m(t(a)) = t^m(t(b)) = t(b)$,
	we have	$t^m(a) = b$.
	By \cref{syncmax:lem:incompressible-stable-pair} (applied with $t^m$, $C$, and $\langle G \cup \{t\} \rangle_S$
	in the roles of $t$, $A$, and~$S$, respectively)
	we find $a \sim_s b$.	
\end{proof}

\begin{theorem} \label{syncmax:thm:unimodal-sync}
	Let $G \le_G \mathcal S_n$.
	Then~$G$ is strongly $\{a,b\}$-primitive
	if and only if it  $G$ synchronizes every unimodal $t \colon [n] \to [n]$ with unimodal pair $a,b$.	
\end{theorem}
\begin{proof}
	First, suppose~$G$ is strongly $\{a,b\}$-primitive.
	Let $t\colon [n] \to [n]$ be unimodal with unimodal pair~$a,b$.
	By \cref{syncmax:lem:unimodal-stable},
	we have $a \sim_s b$ in $\langle G \cup \{t\} \rangle_S$.
	Since $G$ is $\{a,b\}$-primitive, this yield that $\sim_s$
	equals the universal relation. 
	By \cref{prelim:cor:stability-sync} it follows that 
	$\langle G \cup \{t\} \rangle_S$ contains a constant map, i.e., $G$ synchronizes~$t$.
	
	Conversely, each simple idempotent mapping $a$ to~$b$ or vice versa 
	is a unimodal with unimodal pair $a,b$.
	Thus, by \cref{sync_max:cor:2-sets_transitive}, $G$ is strongly $\{a,b\}$-primitive.	
\end{proof}

\begin{corollary} \label{syncmax:cor:unimodal-sync}
	A group on a finite domain of size at least three is primitive
	if and only if it synchronizes every unimodal transformation.	
\end{corollary}
\begin{proof}
	This follows from \cref{sync_max:prop:primitive_Aprimitive,syncmax:thm:unimodal-sync}.
\end{proof}

One-point singulars are unimodal
and with \cref{prelim:lem:transitive-grp-one-point-singular}
for two points mapped to a common state we find
an idempotent one-point singular having 
having these two points as an unimodal pair.
Thus \cref{syncmax:thm:unimodal-sync} also 
implies \cref{syncmax:cor:Aprim-idempt,sync_max:thm:primtiive_kerneltype_k1111}.

\begin{remark}
	Comparing \cref{syncmax:thm:unimodal-sync} and~\cref{syncmax:thm:Aprim-idempt}
	the question arises whether we can characterize
	transitive $A$-primite groups for sets $A$ with more than three points
	by the property that such a group synchronizes every unimodal transformation
	having a unimodal pair in~$A$. We do not know how to obtain such a characterization.	
\end{remark}